\documentclass[10pt,a4paper]{amsart}
\usepackage{amsmath,amssymb,amsthm,graphicx}
\usepackage{verbatim}
\newenvironment{reference}{\comment}{\endcomment}
%\newenvironment{reference}{}{}
\newenvironment{slogan}{\comment}{\endcomment}
\newenvironment{history}{\comment}{\endcomment}

% For commutative diagrams we use Xy-pic
\usepackage[all]{xy}

% We use 2cell for 2-commutative diagrams.
\xyoption{2cell}
% Theorem environments.
%
\theoremstyle{plain}
\newtheorem{theorem}[subsection]{Theorem}
\newtheorem{proposition}[subsection]{Proposition}
\newtheorem{lemma}[subsection]{Lemma}

\theoremstyle{definition}
\newtheorem{definition}[subsection]{Definition}
\newtheorem{example}[subsection]{Example}
\newtheorem{exercise}[subsection]{Exercise}
\newtheorem{situation}[subsection]{Situation}

\theoremstyle{remark}
\newtheorem{remark}[subsection]{Remark}
\newtheorem{remarks}[subsection]{Remarks}

\numberwithin{equation}{subsection}

% Macros
%
\def\lim{\mathop{\mathrm{lim}}\nolimits}
\def\colim{\mathop{\mathrm{colim}}\nolimits}
\def\Spec{\mathop{\mathrm{Spec}}}
\def\Hom{\mathop{\mathrm{Hom}}\nolimits}
\def\Ext{\mathop{\mathrm{Ext}}\nolimits}
\def\SheafHom{\mathop{\mathcal{H}\!\mathit{om}}\nolimits}
\def\SheafExt{\mathop{\mathcal{E}\!\mathit{xt}}\nolimits}
\def\Sch{\mathit{Sch}}
\def\Mor{\mathop{\mathrm{Mor}}\nolimits}
\def\Ob{\mathop{\mathrm{Ob}}\nolimits}
\def\Sh{\mathop{\mathit{Sh}}\nolimits}
\def\NL{\mathop{N\!L}\nolimits}
\def\CH{\mathop{\mathrm{CH}}\nolimits}
\def\proetale{{pro\text{-}\acute{e}tale}}
\def\etale{{\acute{e}tale}}
\def\QCoh{\mathit{QCoh}}
\def\Ker{\mathop{\mathrm{Ker}}}
\def\Im{\mathop{\mathrm{Im}}}
\def\Coker{\mathop{\mathrm{Coker}}}
\def\Coim{\mathop{\mathrm{Coim}}}

% Boxtimes
%
\DeclareMathSymbol{\boxtimes}{\mathbin}{AMSa}{"02}

%
% Macros for moduli stacks/spaces
%
\def\QCohstack{\mathcal{QC}\!\mathit{oh}}
\def\Cohstack{\mathcal{C}\!\mathit{oh}}
\def\Spacesstack{\mathcal{S}\!\mathit{paces}}
\def\Quotfunctor{\mathrm{Quot}}
\def\Hilbfunctor{\mathrm{Hilb}}
\def\Curvesstack{\mathcal{C}\!\mathit{urves}}
\def\Polarizedstack{\mathcal{P}\!\mathit{olarized}}
\def\Complexesstack{\mathcal{C}\!\mathit{omplexes}}
% \Pic is the operator that assigns to X its picard group, usage \Pic(X)
% \Picardstack_{X/B} denotes the Picard stack of X over B
% \Picardfunctor_{X/B} denotes the Picard functor of X over B
\def\Pic{\mathop{\mathrm{Pic}}\nolimits}
\def\Picardstack{\mathcal{P}\!\mathit{ic}}
\def\Picardfunctor{\mathrm{Pic}}
\def\Deformationcategory{\mathcal{D}\!\mathit{ef}}
\usepackage{cleveref}
\crefname{originallemma}{original lemma}{original lemmas}
\Crefname{originallemma}{Original lemma}{Original lemmas}
\makeatletter
\newlabel{coherent-lemma-flat-base-change-cohomology}{{02KH}{source}{}{sourcepointer-1}{}}
\newlabel{coherent-lemma-flat-base-change-cohomology@cref}{{[originallemma][0][]02KH}{[1][0][]source}{}{}{}}
\makeatother
\crefname{originalreference}{original reference}{original references}
\Crefname{originalreference}{Original reference}{Original references}
\makeatletter
\newlabel{lemma-closed-immersion-dim-1-CM}{{0E33}{source}{}{sourcepointer-2}{}}
\newlabel{lemma-closed-immersion-dim-1-CM@cref}{{[originalreference][0][]0E33}{[1][0][]source}{}{}{}}
\makeatother
\crefname{originalreference}{original reference}{original references}
\Crefname{originalreference}{Original reference}{Original references}
\makeatletter
\newlabel{lemma-degree-more-than-2g}{{0E3C}{source}{}{sourcepointer-3}{}}
\newlabel{lemma-degree-more-than-2g@cref}{{[originalreference][0][]0E3C}{[1][0][]source}{}{}{}}
\makeatother
\crefname{originalreference}{original reference}{original references}
\Crefname{originalreference}{Original reference}{Original references}
\makeatletter
\newlabel{lemma-sanity-check-duality}{{0E32}{source}{}{sourcepointer-4}{}}
\newlabel{lemma-sanity-check-duality@cref}{{[originalreference][0][]0E32}{[1][0][]source}{}{}{}}
\makeatother
\crefname{originalreference}{original reference}{original references}
\Crefname{originalreference}{Original reference}{Original references}
\makeatletter
\newlabel{lemma-vanishing-on-gorenstein}{{0E3E}{source}{}{sourcepointer-5}{}}
\newlabel{lemma-vanishing-on-gorenstein@cref}{{[originalreference][0][]0E3E}{[1][0][]source}{}{}{}}
\makeatother
\crefname{originallemma}{original lemma}{original lemmas}
\Crefname{originallemma}{Original lemma}{Original lemmas}
\makeatletter
\newlabel{varieties-lemma-bijection-irreducible-components}{{038F}{source}{}{sourcepointer-6}{}}
\newlabel{varieties-lemma-bijection-irreducible-components@cref}{{[originallemma][0][]038F}{[1][0][]source}{}{}{}}
\makeatother
\crefname{originallemma}{original lemma}{original lemmas}
\Crefname{originallemma}{Original lemma}{Original lemmas}
\makeatletter
\newlabel{varieties-lemma-chi-tensor-finite}{{0AYT}{source}{}{sourcepointer-7}{}}
\newlabel{varieties-lemma-chi-tensor-finite@cref}{{[originallemma][0][]0AYT}{[1][0][]source}{}{}{}}
\makeatother
\crefname{originallemma}{original lemma}{original lemmas}
\Crefname{originallemma}{Original lemma}{Original lemmas}
\makeatletter
\newlabel{varieties-lemma-degree-base-change}{{0B59}{source}{}{sourcepointer-8}{}}
\newlabel{varieties-lemma-degree-base-change@cref}{{[originallemma][0][]0B59}{[1][0][]source}{}{}{}}
\makeatother
\crefname{originallemma}{original lemma}{original lemmas}
\Crefname{originallemma}{Original lemma}{Original lemmas}
\makeatletter
\newlabel{varieties-lemma-geometrically-reduced-any-base-change}{{035Z}{source}{}{sourcepointer-9}{}}
\newlabel{varieties-lemma-geometrically-reduced-any-base-change@cref}{{[originallemma][0][]035Z}{[1][0][]source}{}{}{}}
\makeatother
\crefname{originallemma}{original lemma}{original lemmas}
\Crefname{originallemma}{Original lemma}{Original lemmas}
\makeatletter
\newlabel{varieties-lemma-globally-generated-base-change}{{0B57}{source}{}{sourcepointer-10}{}}
\newlabel{varieties-lemma-globally-generated-base-change@cref}{{[originallemma][0][]0B57}{[1][0][]source}{}{}{}}
\makeatother
\title{Global generation of an invertible sheaf of degree at least 2h\textasciicircum{}1 on every reduced connected closed subscheme of dimension 1 of a proper curve}
\author{The Stacks Project Authors}
\date{stacks-project:curves.tex:local-snapshot-20261012; GNU FDL 1.2 or later}
\begin{document}
\maketitle
\noindent\textit{Editorial note.} All delivered mathematical text is the author's own, copied from the pinned Stacks Project chapter file curves.tex; no statement or proof step was rewritten, reordered or removed.\par\medskip
\noindent\textit{Editorial note.} The counted proof is the author's complete argument: the author's own induction with the purely transcendental field replacement, the ideal-sheaf contradiction setup, the duality dimension count, the generic-point subspace argument, the factorization through Y with both commutative diagrams, and the closing boundary-map contradiction are all retained verbatim.\par\medskip
\noindent\textit{Editorial note.} The label lemma-global-generation names the counted statement, and the proof environment immediately following it in the source is delivered verbatim as its proof.\par\medskip
\noindent\textit{Editorial note.} The counted statement is carried by the author's own lemma environment in the chapter source.\par\medskip
\noindent\textit{Allowed standard prerequisites (editorial).}\begin{itemize}
\item Duality for a proper curve. Same chapter, Serre duality on a proper scheme of dimension 1 over a field: H\textasciicircum{}1(X, F) is dual to Hom(F, omega\_X), used with Lemma lemma-vanishing-on-gorenstein and the deg-chi formula. Supplies the dimension identity dim H\textasciicircum{}0(X, IL) = deg(L) + chi(O\_X) > dim H\textasciicircum{}0(X, omega\_X) that forces the chosen map phi to vanish on global sections.
\item Behaviour of global generation and genus under base field extension. Varieties, Lemmas varieties-lemma-globally-generated-base-change, varieties-lemma-degree-base-change, varieties-lemma-geometrically-reduced-any-base-change, varieties-lemma-bijection-irreducible-components; Cohomology of Schemes, Lemma coherent-lemma-flat-base-change-cohomology. Justifies replacing the ground field k by a purely transcendental extension K, so the union-of-subspaces argument may assume k infinite.
\end{itemize}\medskip
\ifdefined\hypertarget\hypertarget{sourcepointer-1}{}\fi
\noindent\textit{Source reference.} Original lemma 02KH: sources/stacks-chow/tags.txt, source lines 3372--3372. Prerequisite pointer invoked at proof line(s) 5780; it names a standard fact of a sibling chapter that the counted argument uses, and every step of the counted proof built on it is retained verbatim in the delivered body.\par\medskip
\ifdefined\hypertarget\hypertarget{sourcepointer-2}{}\fi
\noindent\textit{Source reference.} Original reference 0E33: sources/stacks-curves/pinned-original-source.tex, source lines 976--976. Prerequisite pointer invoked at proof line(s) 5849; it names a standard fact of the same chapter that the counted argument uses, and every step of the counted proof built on it is retained verbatim in the delivered body.\par\medskip
\ifdefined\hypertarget\hypertarget{sourcepointer-3}{}\fi
\noindent\textit{Source reference.} Original reference 0E3C: sources/stacks-curves/pinned-original-source.tex, source lines 5593--5593. Prerequisite pointer invoked at proof line(s) 5770; it names a standard fact of the same chapter that the counted argument uses, and every step of the counted proof built on it is retained verbatim in the delivered body.\par\medskip
\ifdefined\hypertarget\hypertarget{sourcepointer-4}{}\fi
\noindent\textit{Source reference.} Original reference 0E32: sources/stacks-curves/pinned-original-source.tex, source lines 927--927. Prerequisite pointer invoked at proof line(s) 5781; it names a standard fact of the same chapter that the counted argument uses, and every step of the counted proof built on it is retained verbatim in the delivered body.\par\medskip
\ifdefined\hypertarget\hypertarget{sourcepointer-5}{}\fi
\noindent\textit{Source reference.} Original reference 0E3E: sources/stacks-curves/pinned-original-source.tex, source lines 5680--5680. Prerequisite pointer invoked at proof line(s) 5811; it names a standard fact of the same chapter that the counted argument uses, and every step of the counted proof built on it is retained verbatim in the delivered body.\par\medskip
\ifdefined\hypertarget\hypertarget{sourcepointer-6}{}\fi
\noindent\textit{Source reference.} Original lemma 038F: sources/stacks-chow/tags.txt, source lines 4234--4234. Prerequisite pointer invoked at proof line(s) 5779; it names a standard fact of a sibling chapter that the counted argument uses, and every step of the counted proof built on it is retained verbatim in the delivered body.\par\medskip
\ifdefined\hypertarget\hypertarget{sourcepointer-7}{}\fi
\noindent\textit{Source reference.} Original lemma 0AYT: sources/stacks-chow/tags.txt, source lines 13801--13801. Prerequisite pointer invoked at proof line(s) 5801; it names a standard fact of a sibling chapter that the counted argument uses, and every step of the counted proof built on it is retained verbatim in the delivered body.\par\medskip
\ifdefined\hypertarget\hypertarget{sourcepointer-8}{}\fi
\noindent\textit{Source reference.} Original lemma 0B59: sources/stacks-chow/tags.txt, source lines 14027--14027. Prerequisite pointer invoked at proof line(s) 5777; it names a standard fact of a sibling chapter that the counted argument uses, and every step of the counted proof built on it is retained verbatim in the delivered body.\par\medskip
\ifdefined\hypertarget\hypertarget{sourcepointer-9}{}\fi
\noindent\textit{Source reference.} Original lemma 035Z: sources/stacks-chow/tags.txt, source lines 4146--4146. Prerequisite pointer invoked at proof line(s) 5778; it names a standard fact of a sibling chapter that the counted argument uses, and every step of the counted proof built on it is retained verbatim in the delivered body.\par\medskip
\ifdefined\hypertarget\hypertarget{sourcepointer-10}{}\fi
\noindent\textit{Source reference.} Original lemma 0B57: sources/stacks-chow/tags.txt, source lines 14025--14025. Prerequisite pointer invoked at proof line(s) 5776; it names a standard fact of a sibling chapter that the counted argument uses, and every step of the counted proof built on it is retained verbatim in the delivered body.\par\medskip
\begin{lemma}
\label{lemma-global-generation}
Let $k$ be a field. Let $X$ be a proper scheme over $k$
which is reduced, connected, and of dimension $1$.
Let $\mathcal{L}$ be an invertible $\mathcal{O}_X$-module.
Assume that for every reduced connected closed subscheme
$Y \subset X$ of dimension $1$ we have
$$
\deg(\mathcal{L}|_Y) \geq 2\dim_k H^1(Y, \mathcal{O}_Y)
$$
Then $\mathcal{L}$ is globally generated.
\end{lemma}

\begin{proof}
By induction on the number of irreducible components of $X$.
If $X$ is irreducible, then the lemma holds by
Lemma \ref{lemma-degree-more-than-2g}
applied to $X$ viewed as a scheme over the field
$k' = H^0(X, \mathcal{O}_X)$. Assume $X$ is not irreducible.
Before we continue, if $k$ is finite, then we replace $k$
by a purely transcendental extension $K$. This is allowed by
Varieties, Lemmas
\ref{varieties-lemma-globally-generated-base-change},
\ref{varieties-lemma-degree-base-change},
\ref{varieties-lemma-geometrically-reduced-any-base-change}, and
\ref{varieties-lemma-bijection-irreducible-components},
Cohomology of Schemes, Lemma \ref{coherent-lemma-flat-base-change-cohomology},
Lemma \ref{lemma-sanity-check-duality} and the elementary fact that
$K$ is geometrically integral over $k$.

\medskip\noindent
Assume that $\mathcal{L}$ is not globally generated to get a contradiction.
Then we may choose a coherent ideal sheaf
$\mathcal{I} \subset \mathcal{O}_X$ such that
$H^0(X, \mathcal{I}\mathcal{L}) = H^0(X, \mathcal{L})$
and such that $\mathcal{O}_X/\mathcal{I}$ is nonzero with
support of dimension $0$. For example, take $\mathcal{I}$
the ideal sheaf of any closed point in the common
vanishing locus of the global sections of $\mathcal{L}$.
We consider the short exact sequence
$$
0 \to \mathcal{I}\mathcal{L} \to \mathcal{L} \to 
\mathcal{L}/\mathcal{I}\mathcal{L} \to 0
$$
Since the support of $\mathcal{L}/\mathcal{I}\mathcal{L}$
has dimension $0$ we see that $\mathcal{L}/\mathcal{I}\mathcal{L}$
is generated by global sections
(Varieties, Lemma \ref{varieties-lemma-chi-tensor-finite}).
From the short exact sequence,
and the fact that $H^0(X, \mathcal{I}\mathcal{L}) = H^0(X, \mathcal{L})$
we get an injection
$H^0(X, \mathcal{L}/\mathcal{I}\mathcal{L}) \to H^1(X, \mathcal{I}\mathcal{L})$.

\medskip\noindent
Recall that the $k$-vector space $H^1(X, \mathcal{I}\mathcal{L})$
is dual to $\Hom(\mathcal{I}\mathcal{L}, \omega_X)$.
Choose $\varphi : \mathcal{I}\mathcal{L} \to \omega_X$.
By Lemma \ref{lemma-vanishing-on-gorenstein} we have
$H^1(X, \mathcal{L}) = 0$. Hence
$$
\dim_k H^0(X, \mathcal{I}\mathcal{L}) = \dim_k H^0(X, \mathcal{L}) =
\deg(\mathcal{L}) + \chi(\mathcal{O}_X) > \dim_k H^1(X, \mathcal{O}_X) =
\dim_k H^0(X, \omega_X)
$$
We conclude that $\varphi$ is not injective on global sections, in particular
$\varphi$ is not injective. For every generic point $\eta \in X$
of an irreducible component of $X$ denote
$V_\eta \subset \Hom(\mathcal{I}\mathcal{L}, \omega_X)$ the $k$-subvector
space consisting of those $\varphi$ which are zero at $\eta$.
Since every associated point of $\mathcal{I}\mathcal{L}$
is a generic point of $X$, the above shows that
$\Hom(\mathcal{I}\mathcal{L}, \omega_X) = \bigcup V_\eta$.
As $X$ has finitely many generic points and $k$ is infinite, we conclude
$\Hom(\mathcal{I}\mathcal{L}, \omega_X) = V_\eta$ for some $\eta$.
Let $\eta \in C \subset X$ be the corresponding irreducible component.
Let $Y \subset X$ be the union of the other irreducible components
of $X$. Then $Y$ is a nonempty reduced closed subscheme not equal to $X$.
Let $\mathcal{J} \subset \mathcal{O}_X$ be the ideal sheaf of $Y$.
Please keep in mind that the support of $\mathcal{J}$ is $C$.

\medskip\noindent
Let $\varphi : \mathcal{I}\mathcal{L} \to \omega_X$ be arbitrary.
Since $\mathcal{J}\mathcal{I}\mathcal{L}$
has no embedded associated points
(as a submodule of $\mathcal{L}$) and as $\varphi$ is zero
in the generic point $\eta$ of the support of $\mathcal{J}$, we find that
$\varphi$ factors as
$$
\mathcal{I}\mathcal{L} \to
\mathcal{I}\mathcal{L}/\mathcal{J}\mathcal{I}\mathcal{L} \to \omega_X
$$
We can view $\mathcal{I}\mathcal{L}/\mathcal{J}\mathcal{I}\mathcal{L}$
as the pushforward of a coherent sheaf on $Y$ which by abuse of
notation we indicate with the same symbol.
Since $\omega_Y = \SheafHom(\mathcal{O}_Y, \omega_X)$
by Lemma \ref{lemma-closed-immersion-dim-1-CM}
we find a factorization
$$
\mathcal{I}\mathcal{L} \to
\mathcal{I}\mathcal{L}/
\mathcal{J}\mathcal{I}\mathcal{L} 
\xrightarrow{\varphi'} \omega_Y \to \omega_X
$$
of $\varphi$. Let $\mathcal{I}' \subset \mathcal{O}_Y$ be the
image of $\mathcal{I} \subset \mathcal{O}_X$. There is a surjective map
$\mathcal{I}\mathcal{L}/\mathcal{J}\mathcal{I}\mathcal{L} \to
\mathcal{I}'\mathcal{L}|_Y$ whose kernel is supported in closed points.
Since $\omega_Y$ is a Cohen-Macaulay module on $Y$, the map $\varphi'$
factors through a map
$\varphi'' : \mathcal{I}'\mathcal{L}|_Y \to \omega_Y$.
Thus we have commutative diagrams
$$
\vcenter{
\xymatrix{
0 \ar[r] &
\mathcal{I}\mathcal{L} \ar[r] \ar[d] &
\mathcal{L} \ar[r] \ar[d] &
\mathcal{L}/\mathcal{I}\mathcal{L} \ar[r] \ar[d] &
0 \\
0 \ar[r] &
\mathcal{I}'\mathcal{L}|_Y \ar[r] &
\mathcal{L}|_Y \ar[r] &
\mathcal{L}|_Y/\mathcal{I}'\mathcal{L}|_Y \ar[r] &
0
}
}
\quad\text{and}\quad
\vcenter{
\xymatrix{
\mathcal{I}\mathcal{L} \ar[r]_\varphi \ar[d] & \omega_X \\
\mathcal{I}'\mathcal{L}|_Y \ar[r]^{\varphi''} & \omega_Y \ar[u]
}
}
$$
Now we can finish the proof as follows:
Since for every $\varphi$ we have a $\varphi''$ and since
$\omega_X \in \textit{Coh}(\mathcal{O}_X)$
represents the functor $\mathcal{F} \mapsto \Hom_k(H^1(X, \mathcal{F}), k)$,
we find that
$H^1(X, \mathcal{I}\mathcal{L}) \to H^1(Y, \mathcal{I}'\mathcal{L}|_Y)$
is injective. Since the boundary
$H^0(X, \mathcal{L}/\mathcal{I}\mathcal{L}) \to H^1(X, \mathcal{I}\mathcal{L})$
is injective, we conclude the composition
$$
H^0(X, \mathcal{L}/\mathcal{I}\mathcal{L}) \to
H^0(X, \mathcal{L}|_Y/\mathcal{I}'\mathcal{L}|_Y) \to
H^1(X, \mathcal{I}'\mathcal{L}|_Y)
$$
is injective. Since
$\mathcal{L}/\mathcal{I}\mathcal{L} \to
\mathcal{L}|_Y/\mathcal{I}'\mathcal{L}|_Y$
is a surjective map of coherent modules whose supports have
dimension $0$, we see that the first map
$H^0(X, \mathcal{L}/\mathcal{I}\mathcal{L}) \to
H^0(X, \mathcal{L}|_Y/\mathcal{I}'\mathcal{L}|_Y)$
is surjective (and hence bijective).
But by induction we have that $\mathcal{L}|_Y$ is globally
generated (if $Y$ is disconnected this still works of course)
and hence the boundary map
$$
H^0(X, \mathcal{L}|_Y/\mathcal{I}'\mathcal{L}|_Y) \to
H^1(X, \mathcal{I}'\mathcal{L}|_Y)
$$
cannot be injective.
This contradiction finishes the proof.
\end{proof}

\end{document}
